\chapter*{Préface}
\addcontentsline{toc}{chapter}{Préface}

La classe de Terminale marque l'aboutissement de plusieurs années d'apprentissage des mathématiques et l'entrée dans une épreuve décisive : le Baccalauréat. C'est dans cet esprit qu'\emph{Objectif Bac} a été conçu : offrir aux élèves de Série S un outil complet, structuré et fidèle au programme officiel malgache, pour aborder cette échéance avec méthode et confiance.

Chaque chapitre de cet ouvrage suit la même logique. Un résumé de cours rassemble les définitions, propriétés et théorèmes essentiels, illustrés autant que possible par des schémas et des exemples concrets. Des encadrés « Méthode » détaillent, étape par étape, les démarches à connaître pour résoudre les types de questions les plus fréquemment posés à l'examen. Des exercices résolus permettent ensuite de voir ces méthodes appliquées en détail, avant une série d'exercices gradués — du plus accessible au plus exigeant — pour progresser à son rythme.

Le chapitre \emph{Entraînement}, en fin d'ouvrage, propose une sélection de sujets types conformes à l'esprit du Baccalauréat, couvrant l'ensemble des notions au programme. Ils permettront à l'élève de se confronter à des épreuves complètes, dans les conditions qui seront les siennes le jour de l'examen.

Réussir en mathématiques ne tient pas à un don particulier, mais à la régularité du travail et à la compréhension progressive des raisonnements. Nous espérons que ce manuel accompagnera chaque élève dans cette démarche, jusqu'à la réussite de son Baccalauréat.

\bigskip
\begin{flushright}
Mamiarilala
\end{flushright}
